\documentclass[CMPE]{KGCOEReport}

\usepackage{graphicx}
\usepackage{newtxtext, newtxmath}
\usepackage{parskip}
\usepackage{karnaugh-map}
\usepackage{float}
\usepackage{pdfpages}

\graphicspath{{./assets/}}

\newcommand{\classCode}{CMPE 250}
\newcommand{\name}{Aden Perry}
\newcommand{\LabSectionNum}{L1}
\newcommand{\LabInstructor}{Chongzhou Fang}

\newcommand{\TAs}{Eddie Satowski \\ Khoi Pham}
\newcommand{\LectureSectionNum}{01}
\newcommand{\LectureInstructor}{Chongzhou Fang}
\newcommand{\exerciseNumber}{4}
\newcommand{\exerciseDescription}{Iteration and Subroutines}
\newcommand{\dateDone}{2026-02-12}
\newcommand{\dateSubmitted}{2025-02-21}

\begin{document}

\maketitle

\section*{Abstract}

The purpose of this exercise was to initialize a Universal Asynchronous
Reciever / Transmitter (UART) on an NXP Freedom Board KL05Z, as well as to
design, implement, and test an interactable menu and subroutines to print
characters to and recieve characters from a TTY. The subroutines waited for the
\texttt{TDRE} or \texttt{RDRF} flag to be set, then read or wrote data to or
from the \texttt{UART\_D} register. The interactable menu was implemented by
waiting for a character to be input, and if the character was one of C, N, V,
or Z, branching to a library-provided subroutine to print the menu text. The
program was translated, compiled, and flashed to the KL05Z. It was tested with
various inputs. All inputs worked as expected; the exercise was successful.

\section*{Procedure}

A Keil MDK-ARM project was created. Provided files \texttt{MKL05Z4.s},
\texttt{Start.s}, and \linebreak \texttt{Exercise05\_Lib.lib} were added to the project
directory. Files \texttt{Start.s} and \linebreak \texttt{Exercise05\_Lib.lib} were added
to Source Group 1 of the project.

A subroutine for initializing the UART on the KL05Z to use port B pins 1 and 2
with 8 data bits, no parity, 1 stop bit, and 9600 baud was taken from provided
code and modified so that no registers other than \texttt{LR}, \texttt{PC}, and
\texttt{PSR} were changed after execution.

A subroutine to get a character from the UART data register was written. It
first waited for the \texttt{RDRF} flag to be set, ensuring that data was ready
to be read, then read a byte from \texttt{UART0\_D} into \texttt{R0}.

A subroutine to put a character into the UART data register was written. It
first waited for the \texttt{TDRE} flag to be set, ensuring that data was ready
to be transmitted, than wrote a byte from \texttt{R0} into \texttt{UART0\_D}.

A main program to display an interactive menu was written. It used a library to
display prompt text and responses. The library dependend on the PutChar
subroutine. The main program first polled for input using GetChar, then
converted any lowercase letter to an uppercase letter. The converted inputs
were compared against letters C, N, V, and Z. If any inputs matched, the
library was used to print information about the relevant \texttt{PSR} flag.

The program was translated and built. Target options were modified to allow the
debugger to work with the flashed chip. The program was flashed to the KL05Z.
The windows program PuTTY was configured to communicate with the chip. The
program was tested with a variety of inputs, including A--Z, a--z, and special
characters. Only the characters c, C, n, N, v, V, z, and Z were printed and
brought up the relevant menus. Occasionally, after inputting a long sequence of
characters, the menu stopped responding. It is unclear whether this was caused
due to a fault in the assembly program or beacuse the TTY was misconfigured.

\section*{Results}

The MyCode area of the program had a byte size of 344. The exeucutable area of
of the MyCode area had a byte size of 282. It's memory addresses spanned from
0x0 to 0x11A. The main program spanned from 0x6 to 0x82, with a byte size of
124. The GetChar subroutine spanned from memory address 0xEE to memory address
0x102, with a byte size of 20. The PutChar subroutine spanned from memory
address 0x104 to memory address 0x118, with a byte size of 20. The equal byte
sizes of GetChar and PutChar make sense — the only differences between them is
that they read different flags and do different I/O operations to the
\texttt{UART0\_D} register.

\section*{Conclusion}


The object of this exercise was to initialize a Universal Asynchronous
Reciever / Transmitter (UART) on an NXP Freedom Board KL05Z, as well as to
design, implement, and test an interactable menu and subroutines to print
characters to and recieve characters from a TTY. This was accomplished through
the usage of a template to initialize the program and by adding subroutines to
perform the logic. A library was used to determine the text to be displayed.
GetChar and PutChar read a flag from status register 1 and, depending on the
value, stored / wrote a character to / from \texttt{UART0\_D}. The use of the
UART will helpful in any future work with a chip that needs to communicate with
another device.

\includepdf[pages=-]{assets/prelab}

\end{document}
